% ============================================================
%  paper-export :: GB 预设（中文学位论文 / 中文期刊）
%  配合: -V documentclass=ctexbook --top-level-division=chapter
%
%  注入点: pandoc --include-in-header（位于所有 \usepackage 之后，
%          因此这里的设置可以覆盖 pandoc 默认模板的行为）
% ============================================================

% ---------- 1. 页面：左侧留装订线 ----------
\usepackage{geometry}
\geometry{a4paper,
          top=2.54cm,bottom=2.54cm,left=3.17cm,right=2.54cm,
          headheight=14pt,headsep=0.55cm,footskip=1.05cm}

% 保留 CJK 字符之间的空格（否则「第一章 绪论」会被挤成「第一章绪论」）
\xeCJKsetup{CJKspace=true}

% ---------- 2. 正文：首行缩进 2 字符，1.5 倍行距 ----------
\usepackage{indentfirst}
\AtBeginDocument{\setlength{\parindent}{2em}}
\linespread{1.5}\selectfont
\setlength{\parskip}{0pt plus 1pt}

% ---------- 3. 标题层级 ----------
%   章 → 三号黑体居中，自动编号「第一章」，另起一页（\chapter 自带）
%   节 → 四号黑体
%   条 → 小四黑体
%   款 → 小四黑体
\ctexset{
  chapter/name             = {第,章},
  chapter/number           = \chinese{chapter},
  chapter/format           = \heiti\zihao{3}\centering,
  chapter/beforeskip       = 0pt,
  chapter/afterskip        = 24pt,
  section/format           = \heiti\zihao{4}\raggedright,
  section/beforeskip       = 18pt,
  section/afterskip        = 12pt,
  subsection/format        = \heiti\zihao{-4}\raggedright,
  subsection/beforeskip    = 12pt,
  subsection/afterskip     = 6pt,
  subsubsection/format     = \heiti\zihao{-4}\raggedright,
  subsubsection/beforeskip = 9pt,
  subsubsection/afterskip  = 3pt,
  paragraph/format         = \songti\bfseries\zihao{-4}\raggedright,
  paragraph/runin          = false,
}

% ---------- 4. 目录 ----------
\renewcommand{\contentsname}{目\hspace{1em}录}

\makeatletter
\renewcommand*\l@chapter[2]{%
  \vskip 6pt plus 1pt
  {\heiti\zihao{-4}\@dottedtocline{0}{0em}{2.4em}{#1}{#2}}}
\renewcommand*\l@section[2]{%
  \vskip 2pt plus 1pt
  {\songti\zihao{-4}\@dottedtocline{1}{1.6em}{2.4em}{#1}{#2}}}
\renewcommand*\l@subsection[2]{%
  {\songti\zihao{5}\@dottedtocline{2}{4.0em}{3.0em}{#1}{#2}}}
\renewcommand*\l@subsubsection[2]{%
  {\songti\zihao{5}\@dottedtocline{3}{7.0em}{3.6em}{#1}{#2}}}
\renewcommand{\@pnumwidth}{2.2em}
\renewcommand{\@tocrmarg}{2.8em}

% 目录排完后切回正文页码与页眉样式（正文从阿拉伯数字 1 重新开始）。
%
% 关键在 \aftergroup。pandoc 的 LaTeX 模板把目录整段包在一对花括号里
% （为了让 \hypersetup{linkcolor=toccolor} 只作用于目录）：
%     {\hypersetup{...}\setcounter{tocdepth}{N}\tableofcontents}
% 而 \pagestyle 做的是局部 \let，出了这个分组就弹回 pefront —— 正文页会一直
% 顶着「前置」的页眉、右侧永远没有章节名。偏偏 \pagenumbering 用的是 \gdef，
% 页码看上去完全正常，于是这个 bug 表现为「页码对、页眉不对」，极难察觉。
%
% 用 \globaldefs=1 强行全局化是行不通的：它会把 \protected@write 内部那些
% 本该局部的 \let\protect... 一并全局化，之后写进 .aux 的控制序列会集体丢掉
% 反斜杠，下一轮编译直接报 "Missing \begin{document}"。
% \aftergroup 把切换推迟到分组结束之后、正文开始之前执行，干净且无副作用。
\AtBeginDocument{%
  \let\pe@origtoc\tableofcontents
  \renewcommand{\tableofcontents}{%
    \pe@origtoc
    \aftergroup\pe@beginbody}}
\newcommand{\pe@beginbody}{%
  \clearpage\pagenumbering{arabic}\pagestyle{fancy}}
\makeatother

% ---------- 5. 页眉页脚 ----------
% 学位论文通行做法：页眉居中当前章名 + 横线，页脚居中页码。
% 三节：封面（不计页码）→ 前置（罗马）→ 正文（阿拉伯从 1）
\usepackage{fancyhdr}
\usepackage{lastpage}

\makeatletter
\providecommand\pe@title{}
\AtBeginDocument{\global\let\pe@title\@title}
% 与 report 同理：fancyhdr 切换页面样式时会重置 \chaptermark，
% 因此把它包成宏、在每个 \fancypagestyle 体内重新施加。
\newcommand{\pe@setmarks}{%
  \renewcommand{\chaptermark}[1]{\markboth{\CTEXthechapter\hspace{0.5em}##1}{}}}

\providecommand{\pefootertotal}{0}
\newcommand{\pe@pagenum}{%
  \ifnum\pefootertotal=1\relax
    第~\thepage~页\quad 共~\pageref*{LastPage}~页%
  \else
    \thepage
  \fi}
\newcommand{\pe@hf}{\songti\zihao{-5}\color{pblack}}

\fancypagestyle{fancy}{%
  \pe@setmarks
  \fancyhf{}%
  \fancyhead[C]{\pe@hf\leftmark}%
  \fancyfoot[C]{\pe@hf\pe@pagenum}%
  \renewcommand{\headrulewidth}{0.6pt}%
  \renewcommand{\footrulewidth}{0pt}}
\fancypagestyle{pefront}{%
  \fancyhf{}%
  \fancyhead[C]{\pe@hf\pe@title}%
  \fancyfoot[C]{\pe@hf\thepage}%
  \renewcommand{\headrulewidth}{0.6pt}%
  \renewcommand{\footrulewidth}{0pt}}
\fancypagestyle{empty}{\fancyhf{}%
  \renewcommand{\headrulewidth}{0pt}\renewcommand{\footrulewidth}{0pt}}
% \chapter 会把当页切成 plain：保留页码，去掉页眉
\fancypagestyle{plain}{%
  \fancyhf{}%
  \fancyfoot[C]{\pe@hf\pe@pagenum}%
  \renewcommand{\headrulewidth}{0pt}\renewcommand{\footrulewidth}{0pt}}
\makeatother

% \headrule 必须自己判 \headrulewidth —— 直接画死线会让 \fancypagestyle{empty}
% 里的 headrulewidth=0pt 失效，封面顶端就会漏出一条页眉线。
\renewcommand{\headrule}{%
  \ifdim\headrulewidth>0pt\relax
    {\color{prule}\hrule height \headrulewidth}%
  \fi}
\pagestyle{empty}

% ---------- 6. 图表题注：黑体五号居中，标签与内容间空一格 ----------
\usepackage{caption}
\DeclareCaptionFont{pcap}{\heiti\zihao{5}\color{pblack}}
\captionsetup{font=pcap,labelsep=quad,skip=6pt,justification=centering,
              singlelinecheck=true}
\captionsetup[figure]{position=bottom}
\captionsetup[table]{position=top}

% ---------- 7. 编号按章：图 1-1 / 表 2-3 / 式 (4-5) ----------
\renewcommand{\thefigure}{\thechapter-\arabic{figure}}
\renewcommand{\thetable}{\thechapter-\arabic{table}}
\renewcommand{\theequation}{\thechapter-\arabic{equation}}
\makeatletter
\@addtoreset{figure}{chapter}
\@addtoreset{table}{chapter}
\@addtoreset{equation}{chapter}
\makeatother

